\section{Results}

\subsection{A frozen action-contrast transfer protocol}

The final experiment evaluates target offline adaptation rather than zero-shot
policy transfer.  Source counterfactuals and target counterfactuals are both
generated from matched restored SUMO states, but target roles are separated by
seed (\Cref{fig:method-schematic}).  Seeds 5057 and 6067 are the only target
adaptation seeds.  Seed 8171 is collected after the joint method freeze and is
used once for offline action-regret confirmation.  Seeds 8081 and 9091 are used
only for closed-loop policy rollouts.  Out-of-fold models select the blend but
are never deployed; after selection, each external city receives one refit on
the union of all its valid adaptation groups.

\begin{figure}[!htbp]
\centering
\includegraphics[width=\textwidth]{tsc_anchored_transfer_schematic.pdf}
\caption{Transferred object and frozen evaluation protocol. \textbf{a}, The
method lineage contrasts absolute dense-residual prediction, factored MC-WM
dynamics and the final matched-ranking estimand.  The same-protocol dense
baseline in panel \textbf{b} is deliberately strengthened with the matched
contrast data and target refit budget used by CFCMT. \textbf{b}, A rigid
parent-restricted anchor and an antisymmetric all-pair correction are combined
by a seed-blocked target selector.  The pressure reference defines the contrast
origin; MaxPressure and phase pressure are independent comparators, not
fallback actions.  Target counterfactual labels make the final setting target
offline adaptation rather than zero-shot transfer.}
\label{fig:method-schematic}
\end{figure}

The source bank contains 18 networks in seven city groups: Atlanta, Cologne,
Hangzhou, Ingolstadt, New York, two RESCO synthetic grids and Salt Lake City.
Together they contain 317 controlled traffic lights and 6,076 non-internal
lanes.  The external set contains a four-signal Los Angeles corridor and a
12-signal Jinan grid evaluated under three demand files
(\Cref{fig:network-plate}; detailed counts in \Cref{tab:network-scope}).
City groups, rather than network files, are the development transfer unit, so
another Cologne, Hangzhou or Ingolstadt subnetwork cannot enter the training
fold when that city is held out.

\begin{figure}[!htbp]
\centering
\includegraphics[width=\textwidth]{traffic_signal_transfer_networks.pdf}
\caption{Network shift in development and external confirmation.  Lines show
non-internal SUMO lanes and points show controlled traffic lights.  Blue panels
are the 18 source networks grouped into seven development cities; green panels
are external target geometries.  The Jinan geometry is paired with three demand
files.  The plate documents topology and action-set variation; it does not
imply that the converted networks are calibrated city-wide digital twins.}
\label{fig:network-plate}
\end{figure}

\subsection{Seed-held-out development selected a non-anchor correction}

The method was not accepted because it fit its development seed.  Models and
the city-level selector were fitted using source seeds 2027 and 3037, then
evaluated on seed 4047.  In leave-one-city-out evaluation, the rigid anchor had
city-macro normalized action regret 0.2414 and the selected anchored model had
regret 0.2177, a 9.84\% reduction (\Cref{tab:development-confirmation}).  Six
of seven held-out city groups improved.  Hangzhou was the sole regression,
with an absolute increase of 0.00388, below the frozen 0.01 tolerance.

\begin{table}[t]
\caption{Seed-held-out development confirmation.  Regret is candidate cost
minus the matched-group optimum, normalized by the within-group action range;
lower is better.}
\label{tab:development-confirmation}
\centering
\small
\begin{tabular}{lrrr}
\toprule
Development result & Anchor regret & Selected regret & Relative gain \\
\midrule
Seven-city leave-one-city-out & 0.241 & 0.218 & 9.84\% \\
\bottomrule
\end{tabular}

\end{table}

Earlier variants failed their gates.  A global pairwise correction regressed
on the added Salt Lake domain, the first anchored selector exceeded its
worst-city tolerance, and a cross-fitted selector did not achieve the required
city-level stability.  These failures were retained and the successful
seed-held-out protocol was frozen before external labels were generated
(\Cref{tab:development-history}).

\subsection{External offline scores favoured source-plus-target fitting}

The external selector used every valid adaptation group: 165 in Los Angeles
and 911 across the three Jinan demands.  Holding out each complete adaptation
seed produced two selection folds per city.  Los Angeles selected a local
blend with disagreement cap 1 and confidence multiplier 0.25; Jinan selected a
constant correction weight of 0.9.  The deployment artifacts were then refit
once on all adaptation groups and frozen before seed 8171 was collected.

On seed 8171, Los Angeles regret decreased from 0.1482 to 0.1450 (2.18\%) and
optimal-action accuracy increased from 0.785 to 0.810.  Jinan regret decreased
from 0.4361 to 0.3921 (10.10\%) and optimal-action accuracy increased from
0.296 to 0.342 (\Cref{tab:external-offline,fig:external-confirmation}a).
Equal weighting of the two cities gives a regret reduction from 0.2922 to
0.2685, or 8.09\%, and both city point estimates improve.

The freeze also contained a target-label-only comparator from the same causal
action-advantage family.  It used the same feature semantics and target
adaptation groups, but its target prediction weight was numerically one, so
source-labelled groups did not contribute to its deployed score.  Its
seed-8171 regret was 0.2172 in Los Angeles and 0.4036 in Jinan.  CFCMT therefore
reduced the equal-city point estimate from 0.3104 to 0.2685 (13.48\%).
The paired action-group improvement was 0.07219 in Los Angeles (95\% interval
$[0.01671,0.13512]$, two-sided bootstrap probability value 0.011) and 0.01150
in Jinan (interval $[-0.02250,0.04524]$, probability value 0.5196).  Offline
source contribution was therefore supported in Los Angeles but inconclusive in
Jinan, despite both point estimates favouring CFCMT.

\begin{table}[t]
\caption{External seed-8171 offline confirmation after city-specific method
freeze.  Target only uses the same causal action-ranking family and target
groups but zero deployed source weight.  Adaptation groups are target
counterfactuals from seeds 5057 and 6067; held-out groups are never used for
fitting or selection.  Gains in the table compare CFCMT with target only.}
\label{tab:external-offline}
\centering
\scriptsize
\resizebox{\textwidth}{!}{\begin{tabular}{lrrrrrrr}
\toprule
City & Adapt. groups & Held-out groups & Anchor & Target only & CFCMT & CFCMT vs target only & Selected blend \\
\midrule
Los Angeles & 165 & 79 & 0.148 & 0.217 & 0.145 & 33.24\% & local\_cap\_1\_z\_0p25 \\
Jinan & 911 & 457 & 0.436 & 0.404 & 0.392 & 2.85\% & constant\_alpha\_0p9 \\
\midrule
External macro & -- & 2 cities & 0.292 & 0.310 & 0.269 & 13.48\% & city-specific \\
\bottomrule
\end{tabular}
}
\end{table}

The uncertainty differs sharply by city.  A 10,000-replicate paired group
bootstrap gives a Los Angeles improvement interval of
$[-0.0675,0.0729]$ with a two-sided bootstrap probability value of 0.9396,
whereas Jinan gives $[0.00469,0.08293]$ and 0.0274.  The Los Angeles result is
therefore directional, not independently conclusive.  The external gate is a
predeclared engineering confirmation (at least 2\% city-macro improvement,
both cities improving and no city regression), not a population-level test
over cities.

\begin{figure}[!htbp]
\centering
\includegraphics[width=\textwidth]{tsc_external_confirmation.pdf}
\caption{Independent offline and closed-loop confirmation. \textbf{a},
seed-held-out normalized action regret for development and external targets.
Green points show the selected CFCMT score and grey points the rigid anchor.
\textbf{b}, all-departed waiting time by policy in Los Angeles (circles) and
Jinan (diamonds); bars show the equal-city macro. \textbf{c}, paired CFCMT
waiting-time gains relative to the non-fixed comparators; fixed time remains
visible in panel \textbf{b} and \Cref{tab:external-closed-loop}.  Positive values
favour CFCMT.  Intervals are descriptive paired hierarchical bootstrap
intervals because there are only two independent target cities.}
\label{fig:external-confirmation}
\end{figure}

\subsection{Closed-loop transfer improved dense residual control}

The formal closed-loop matrix contains four target scenarios, two seeds and
seven policies, yielding 56 completed 3,600-s rollouts.  The primary metric is
mean waiting over every departed vehicle.  SUMO \texttt{write-unfinished}
records accrue waiting to the fixed horizon, so unfinished vehicles remain in
the population.  Completed-only waiting, system vehicle-hours and completion
are reported beside it to expose censoring and throughput trade-offs.

Anchored CFCMT achieved 149.17 s all-departed waiting, compared with 159.15 s
for the H2O$^+$-style dense residual, 156.71 s for simulator-only control and
152.91 s for the rigid anchor (\Cref{tab:external-closed-loop}).  Relative to
the dense residual, the reduction was 9.98 s or 6.27\%; both Los Angeles
(8.46 s) and Jinan (11.50 s) improved.  The descriptive paired hierarchical
bootstrap interval was $[2.01,17.19]$ s and 98.43\% of replicates favoured
CFCMT.  Completed-only waiting was also lower (133.06 versus 147.47 s), system
vehicle-hours decreased from 618.15 to 594.14, and completion increased from
0.870 to 0.889.

\begin{table}[t]
\caption{External 3,600-s closed-loop city-macro results.  Scenarios are
weighted equally within each city and the two cities are weighted equally.
Waiting times are seconds; system burden is vehicle-hours.  ``Incidents'' is
the equal-city mean of unique SUMO junction-collision incidents.}
\label{tab:external-closed-loop}
\centering
\scriptsize
\resizebox{\textwidth}{!}{\begin{tabular}{lrrrrrr}
\toprule
Policy & All-departed wait & Completed wait & System veh-h & Unfinished & Completion & Incidents \\
\midrule
Fixed time & 338.58 & 308.72 & 885.18 & 0.133 & 0.661 & 2.00 \\
Dense residual & 159.15 & 147.47 & 618.15 & 0.066 & 0.870 & 0.25 \\
Simulator only & 156.71 & 138.68 & 619.66 & 0.069 & 0.877 & 0.00 \\
Rigid anchor & 152.91 & 138.40 & 606.30 & 0.072 & 0.882 & 0.50 \\
\textbf{Anchored CFCMT} & 149.17 & 133.06 & 594.14 & 0.067 & 0.889 & 0.00 \\
Phase pressure & 143.62 & 125.32 & 583.05 & 0.068 & 0.898 & 0.25 \\
MaxPressure & 141.74 & 135.42 & 580.29 & 0.065 & 0.896 & 0.25 \\
\bottomrule
\end{tabular}
}
\end{table}

The point-estimate gains over simulator-only and rigid-anchor control were
7.54 s (4.81\%) and 3.75 s (2.45\%), respectively, but their intervals crossed
zero (\Cref{tab:external-comparisons}).  Thus the closed-loop evidence is
strongest against dense residual transfer.  Improvement over the anchor is
supported clearly by offline action regret and only directionally by these two
closed-loop cities; the result does not justify attributing every policy gain
to the pairwise correction.

\begin{table}[t]
\caption{Paired CFCMT policy contrasts for all-departed waiting.  Positive
gains favour CFCMT.  Intervals are descriptive, not population-level city
inference.}
\label{tab:external-comparisons}
\centering
\scriptsize
\resizebox{\textwidth}{!}{\begin{tabular}{lrrrrr}
\toprule
Baseline & CFCMT gain (s) & Relative gain & Descriptive 95\% interval & LA gain & Jinan gain \\
\midrule
Fixed time & 189.42 & 55.94\% & [67.11, 308.32] & 302.90 & 75.93 \\
Dense residual & 9.98 & 6.27\% & [2.01, 17.19] & 8.46 & 11.50 \\
Simulator only & 7.54 & 4.81\% & [-3.62, 25.27] & 14.75 & 0.34 \\
Rigid anchor & 3.75 & 2.45\% & [-8.75, 15.70] & 2.36 & 5.14 \\
Phase pressure & -5.55 & -3.86\% & [-9.54, -2.45] & -6.91 & -4.19 \\
MaxPressure & -7.43 & -5.24\% & [-15.17, -1.56] & -12.53 & -2.33 \\
\bottomrule
\end{tabular}
}
\end{table}

\subsection{Fresh target-only comparison exposed network-dependent negative transfer}

The seven-policy v43 matrix did not deploy its frozen target-label-only
baseline.  We first added that comparator to the same eight scenario--seed
cells after seeing the main outcomes.  This post-hoc v90 ablation left both
models unchanged: CFCMT had 149.17 s city-macro waiting and target only had
149.04 s, a 0.12-s disadvantage.  Los Angeles favoured target only by 2.45 s,
whereas Jinan favoured CFCMT by 2.20 s.  Because the comparator was added after
the original result, this matrix is descriptive rather than confirmatory.

We then froze a v91 protocol before generating 64 further simulator seeds,
disjoint from all v43, v89 and v90 roles.  Each seed paired CFCMT and target
only on the Los Angeles scenario and all three Jinan demands, giving 256 paired
cells and 512 completed rollouts.  The primary relative difference averaged
scenarios within city, weighted Los Angeles and Jinan equally within seed and
then averaged seeds; negative values favour CFCMT.

Target only achieved 149.09 s and CFCMT 155.14 s.  The equal-seed relative
difference was $+4.254\%$ (paired 95\% interval $[-0.775,13.211]\%$), and only
29/64 seeds improved (one-sided exact sign-test probability 0.809).  The city
effects had opposite signs: CFCMT improved Jinan by 3.045\% but regressed Los
Angeles by 7.865\%.  The worst seed regressed by 267.56\%, collision incidents
were 25 for CFCMT and 15 for target only, and both had zero teleports.  The
predeclared mean, interval, both-city, worst-seed, improved-fraction, sign-test
and collision gates therefore did not jointly pass, and the fresh source-
contribution confirmation was rejected (\Cref{tab:source-contribution}).

\begin{table}[t]
\caption{Closed-loop source-contribution ablation.  Both policies use identical
target adaptation groups, features, action executor and target scenarios;
target only removes source-labelled samples from deployed prediction.  Positive
$\Delta$ denotes higher waiting for CFCMT.  V90 is post-hoc; only V91 is a
fresh-seed preregistered confirmation.  Incident entries are CFCMT/target only.}
\label{tab:source-contribution}
\centering
\scriptsize
\resizebox{\textwidth}{!}{\begin{tabular}{lrrrrrrrrr}
\toprule
Evidence & Seeds & Paired cells & Target-only wait & CFCMT wait & $\Delta$ wait & LA $\Delta$ & Jinan $\Delta$ & Improved & Incidents C/T \\
\midrule
V90 post-hoc & 2 & 8 & 149.04 & 149.17 & +0.082\% & +1.212\% & -2.292\% & 7/8 & 0/1 \\
V91 fresh (rejected) & 64 & 256 & 149.09 & 155.14 & +4.254\% & +7.865\% & -3.045\% & 29/64 & 25/15 \\
\bottomrule
\end{tabular}
}
\end{table}
\FloatBarrier

\subsection{Classical pressure remained the performance ceiling}

MaxPressure and phase pressure achieved 141.74 and 143.62 s waiting,
respectively, outperforming CFCMT by 7.43 and 5.55 s.  The direction was the
same in both external cities and the descriptive intervals excluded zero.
Their system vehicle-hours and completion ratios were also better than
CFCMT.  This is a substantive boundary, not an implementation detail: the
final CFCMT controller does not fall back to either rule.  The pressure action
used to construct a contrast is merely the zero point of a learned ranking,
and any feasible phase can be selected at deployment.

All 56 rollouts completed with zero teleports.  Across all policies there were
171 raw collision events consolidated into 13 unique junction incidents in
six rollouts.  CFCMT had zero incidents and never recorded more incidents than
its paired comparator.  Because incidents are simulator events and the
matrix contains only eight CFCMT rollouts, this audit is a validity and failure
report rather than a safety-certification claim.

\subsection{A fresh-seed successor stress test did not close the pressure gap}

After v43 was frozen, we separately tested whether a longer-horizon latent
originator and a globally rate-limited executor could turn the pointwise
ranking model into a controller that exceeded phase pressure.  This successor
track is not the anchored all-pair v43 controller and is not pooled with its
two-city result.  It uses a target-state-conditioned 120-s action originator,
one permitted network-wide override at a time and a 120-s global switch
cooldown on the Jinan target.

The legacy global executor appeared favourable on 54 disclosed adaptive seeds
($-0.429\%$ mean relative waiting difference; paired 95\% interval
$[-0.818,-0.033]\%$).  Allowing unlimited cooldown stays instead regressed by
$0.104\%$, and a one-stay budget produced an inconclusive $-0.242\%$ point
estimate; both variants were rejected.  We therefore froze the unchanged
legacy executor before generating 64 new seeds disjoint from all adaptive
seeds and from eight still-sealed prospective seeds.

On those 64 fresh seeds, the method--phase-pressure difference shrank to
$-0.043\%$ (paired 95\% interval $[-0.449,0.370]\%$).  The method improved
34/64 seeds, the one-sided exact sign-test probability was 0.354 and the worst
seed regressed by 5.21\%, just beyond the preregistered 5\% limit.  Teleports
were 0/0 and collision incidents 1/1, but the bootstrap, sign-test and
worst-seed gates failed.  The joint confirmation was therefore rejected
(\Cref{tab:postfreeze-stress}).  This untouched-seed result supports neither
same-network superiority over phase pressure nor a cross-network claim for the
successor executor.

\subsection{Native reinforcement-learning results are protocol checks}

Official RESCO and patched LibSignal experiments confirm that standard RL
baselines can be executed in the project.  The RESCO adapter completed 42
supported full-budget runs, and the LibSignal adapter completed 36 runs across
DQN, PressLight, FRAP and MPLight.  They use their native observations,
rewards, action intervals, training loops and metrics, and therefore cannot be
inserted into the 56-rollout table as if they were paired policy rows.  Their
results and upstream exclusions are reported in
\Cref{tab:official-resco-rl,tab:libsignal-rl}; the same-protocol conclusion is
based only on the external CFCMT matrix above.

\FloatBarrier
